\section{Findings}

\subsection{The field's CIFAR-10 anchor is not a CIFAR-10 result}
Progress in Tsetlin image classification is commonly measured against a convolutional Tsetlin machine
result of \textbf{60.7\%}. Tracing it: it originates in an MSc thesis (University of Agder, 2020)
which, in the words of the paper that cites it, ``utilized a \emph{modified} CIFAR-10 dataset''. It
is cited at third hand; its clause budget, window, $T$, $s$ and epoch count are unknown. The 2019
Convolutional Tsetlin Machine paper reports \textbf{no CIFAR-10 number at all}.

The defensible bars are: best single model \textbf{75.4\%} (64\,000 clauses per class, 250 epochs);
previous single-model state of the art \textbf{75.1\%}; a vanilla convolutional machine
\textbf{69.3\%}; and \textbf{82.8\%} for a \emph{composite} of twenty-two specialists, which is an
ensemble and is not comparable to a single model.

This matters beyond bookkeeping: a programme aiming at 60.7\% and one aiming at 75.4\% are different
programmes, and the gap to CNNs is 14.7 points smaller than the anchor implies.

\subsection{Our harness reproduces the literature, and clears it}
At 20\,000 clauses, 5$\times$5 colour thermometers, $T=3000$, $s=5.0$, budget 32, \textbf{50} epochs:
\textbf{\PreflightVal\% validation, \PreflightTest\% test}, against the source table's
\textbf{64.5\%} at \textbf{250} epochs. The reproduction target was cleared at one fifth of the
training budget. This is the first time in this line of work that a published Tsetlin CIFAR-10 cell
has been matched here.

\subsection{The clause-size budget is worth nine points, not 0.06}
\label{sec:budget}
A clause-size budget caps how many literals a clause may include. Its authors measured it at
\textbf{+0.06} points on CIFAR-2 and \textbf{+0.05} on MNIST-with-convolution, and presented it as an
interpretability and hardware mechanism. It is set to 32 in every reference script of the
state-of-the-art toolbox and appears in none of that paper's tables.

On convolutional CIFAR-10 we measure it at \textbf{+\BudDelta\ points}
(\BudCap\ $\pm$ \BudCapSd\ against \BudUnc\ $\pm$ \BudUncSd, three seeds each, seed intervals
disjoint), at \textbf{4.8$\times$ fewer included literals per image} and \textbf{+1.7\% wall-clock}.
For scale, the largest published single-model CIFAR-10 mechanism, Drop Clause, is worth +5.8.

\pgfplotstabletypeset[
  col sep=comma, string type,
  columns/s/.style={column name=$s$},
  columns/mode/.style={column name=feedback},
  columns/unconstrained/.style={column name=unconstrained},
  columns/capped/.style={column name=budget 32},
  columns/delta/.style={column name=$\Delta$},
  columns/unc_sd/.style={column name={sd}},
  columns/cap_sd/.style={column name={sd}},
  every head row/.style={before row=\toprule, after row=\midrule},
  every last row/.style={after row=\bottomrule},
]{data/budget.csv}

Three things make this a finding about Tsetlin machines rather than about our implementation:

\begin{itemize}
\item \textbf{It survives the exact algorithm.} Under sequential feedback --- the per-example
      algorithm, not our batched approximation --- the effect is \textbf{+\SeqDelta}. The two modes
      differ by 0.71 points, inside the seed band, so batch aggregation is not producing it.
\item \textbf{It is not the specificity parameter in disguise.} This was pre-registered as a rival
      explanation by the team member whose result it threatened, because $s$ also shortens clauses
      and this repository's own prior notes recommended using $s$ instead. Tested: $s{=}2$
      unconstrained reaches \StwoUnc, against \BudCap\ for budget 32 --- a gap of 11.15 points
      against a \SeedBand\ point band. The budget pays at every $s$ tested (+\StwoDelta\ at $s{=}2$,
      +\BudDelta\ at $s{=}10$), and the best cell in the family is $s{=}2$ \emph{with} the budget,
      at \StwoCap. The two dials are complementary.
\item \textbf{It holds at scale.} At 20\,000 clauses, a different encoding and a different $s$, the
      budget is still worth +6.43 points.
\end{itemize}

\subsection{Why: inclusion is a ratchet}
The mechanism was derived from the feedback rules and then measured. Adding a literal is effectively
deterministic; removing one is not. An included literal can only be evicted when the clause
\emph{fails to fire}, at rate $1/s$; and the admission rule asks only whether a literal is
\emph{frequent inside the clause's own match set}, never whether it is \emph{useful}. The result is a
ratchet.

Measured on an unconstrained clause over 30 epochs: clause length \textbf{86.7 $\to$ 170.6}
(+97\%, i.e.\ 68.7\% of all includable literals) while the firing rate moves \textbf{$-$1.4\%} and
validation accuracy \textbf{+1.1} points. Epochs 5--30 alone add \textbf{34.4 literals for +0.32
points}, inside the seed band. \textbf{Clause length had not converged at 30 epochs; accuracy
converged at epoch 4.} Per-literal selectivity is \textbf{4.6$\times$} higher under the budget.

At 20\,000 clauses the ratchet is not merely wasteful but harmful: the unconstrained arm's validation
accuracy \textbf{peaks at epoch 13 and declines by 2.30 points} by epoch 50 while its clauses grow to
188 literals. The machine trains itself backwards for 37 epochs. Only validation-selected reporting
sees this; a last-epoch or last-25-epoch statistic --- what every Tsetlin image paper in the
bibliography uses --- records the degraded number and never sees the peak.

One quantitative prediction of the theory was confirmed and one falsified. The predicted equilibrium
firing rate $f^{*}=1/(1+s)$ holds to \textbf{11--13\%} across a sevenfold movement in the predicted
quantity. But clause length was predicted to fall below 32 at $s{=}2$ and instead stayed at 160: the
admission threshold is \emph{not} what governs clause length. What does govern it is open.
